%------------------------
% One-page resume -- revised with reviewer feedback
%------------------------

\documentclass[letterpaper,10pt]{article}

\usepackage[T1]{fontenc}
\usepackage[
left=0.45in,
right=0.45in,
top=0.37in,
bottom=0.1in
]{geometry}
\usepackage{titlesec}
\usepackage{xcolor}
\usepackage{enumitem}
\usepackage[hidelinks]{hyperref}
\usepackage{fancyhdr}
\usepackage{newtxtext,newtxmath}
\usepackage{microtype}

\pagestyle{fancy}
\fancyhf{}
\fancyfoot{}
\renewcommand{\headrulewidth}{0pt}
\renewcommand{\footrulewidth}{0pt}
\setlength{\footskip}{7pt}

\urlstyle{rm}
\raggedright
\setlength{\tabcolsep}{0pt}
\setlength{\parindent}{0pt}
\setlength{\parskip}{0pt}

% --------- SECTION FORMATTING ---------
\titleformat{\section}
{\fontsize{14}{14.5}\selectfont\scshape\raggedright\bfseries}
{}{0em}{}
[\color{black}\titlerule \vspace{-3pt}]

\titlespacing*{\section}{0pt}{4pt}{3pt}

% --------- EXPERIENCE / RESEARCH HEADERS ---------
\newcommand{\resumeSubheading}[4]{
\vspace{2pt}
\begin{tabular*}{\textwidth}{@{}l@{\extracolsep{\fill}}r@{}}
\textbf{#1} & #2 \\
\textit{#3} & \textit{#4} \\
\end{tabular*}\vspace{-5pt}
}

% --------- PROJECT HEADERS ---------
\newcommand{\resumeProjectHeading}[2]{
\vspace{2pt}
\begin{tabular*}{\textwidth}{@{}l@{\extracolsep{\fill}}r@{}}
\textbf{#1} & \textit{#2} \\[-3pt]
\end{tabular*}
\vspace{-11pt}
}

% --------- PROJECT BULLETS ---------
\newcommand{\resumeProjectItemListStart}{
\begin{itemize}[
leftmargin=0.17in,
itemsep=0pt,
topsep=-2pt,
parsep=0pt,
partopsep=0pt
]
}

\newcommand{\resumeProjectItemListEnd}{
\end{itemize}\vspace{0pt}
}

% --------- EXPERIENCE / RESEARCH BULLETS ---------
\newcommand{\resumeItemListStart}{
\begin{itemize}[
leftmargin=0.17in,
label=\textbullet,
itemsep=0pt,
topsep=0pt,
parsep=0pt,
partopsep=0pt
]
}

\newcommand{\resumeItemListEnd}{
\end{itemize}\vspace{-1pt}
}

%-----------------------------
\begin{document}

%----------HEADING-----------------
\begin{center}
{\LARGE \textbf{\scshape Shivani Swaminathan}} \\[-1pt]
\small
678-800-6094 $|$
\href{mailto:sswaminathan67@gatech.edu}
{\underline{sswaminathan67@gatech.edu}} $|$
\href{https://www.linkedin.com/in/shivani-swaminathan-gatech}
{\underline{linkedin.com/in/shivani}} $|$
\underline{US Citizen}
\end{center}

\vspace{-7pt}

%-----------EDUCATION-----------------
\section{Education}

\begin{tabular*}{\textwidth}{@{}l@{\extracolsep{\fill}}r@{}}
\textbf{Georgia Institute of Technology} & Atlanta, GA \\
\textit{Bachelor of Science / Master of Science in Aerospace Engineering} &
\textit{B.S. Fall 2027 $|$ M.S. Fall 2028} \\
\end{tabular*}

\vspace{0pt}

\textit{GPA: 3.78/4.0 $|$ Dean's List $|$ Relevant Coursework: Aerodynamics, Structural Analysis, Fluids, Thermodynamics, Mechanics of Materials}

\vspace{2pt}

%-----------RESEARCH-----------------
\section{Research}

\resumeSubheading
{Ben T. Zinn Combustion Laboratory}{Atlanta, GA}
{Undergraduate Research Assistant}{May 2026 -- Present}
\resumeItemListStart
\item Improved combustor-model fidelity by 18\% across 8 steady and transient ANSYS Fluent cases by refining 19M-cell meshes and configuring reacting-flow physics for higher-confidence combustion predictions.
\item Validated 3 combustor cases by analyzing temperature, pressure, velocity, and species results with ANSYS engineers, leading to 4 model and setup corrections and a more reliable simulation baseline.
\resumeItemListEnd

\resumeSubheading
{GT Aerospace Systems Design Laboratory (ASDL)}{Atlanta, GA}
{Undergraduate Research Assistant -- Systems Modeling}{August 2025 -- Present}
\resumeItemListStart
\item Eliminated 2 infeasible Mars rover concepts after SysML models exposed violations in mass, power, torque, and climb requirements across 5+ subsystems, preventing further analysis of nonviable designs.
\item Reduced the rover design trade space by 30\% through 30+ ModelCenter sensitivity runs on slope, wheel--soil interaction, traction, and power margins, focusing further analysis on feasible concepts.
\resumeItemListEnd

\resumeSubheading
{GT Aerospace Systems Design Laboratory (ASDL)}{Atlanta, GA}
{Undergraduate Research Assistant -- AI/ML}{June 2026 -- Present}
\resumeItemListStart
\item Developed a Python AI agent that expanded a lunar-base simulation to 12+ autonomous construction tasks, automating landing-pad sequencing, equipment allocation, and task dependencies for repeatable mission-plan evaluation.
\item Integrated 4 lunar infrastructure models across 4 mission scenarios by defining dependencies among landing-pad construction, in-situ resource utilization (ISRU) mining, and power operations, enabling end-to-end surface construction simulation.
\resumeItemListEnd

\resumeSubheading
{UGA Department of Physics and Astronomy}{Athens, GA}
{Undergraduate Research Assistant}{August 2024 -- May 2025}
\resumeItemListStart
\item Analyzed 200+ stellar-flare light curves in Python and modeled activity decay across 3 spectral classes in MATLAB, reducing processing time by 36\% and earning an undergraduate research award.
\resumeItemListEnd

%-----------EXPERIENCE-----------------
\section{Experience}

\resumeSubheading
{GT HyTech Racing}{Atlanta, GA}
{Aerodynamics Engineer}{August 2025 -- Present}
\resumeItemListStart
\item Validated carbon-fiber wing assemblies across 6+ ride-height and angle-of-attack cases in ANSYS Fluent, identifying up to 18\% downforce and 11\% drag variation to select competition wing ride height and angle of attack.
\item Reduced analysis turnaround by 36\% by automating ANSYS Fluent post-processing in Python, enabling same-day comparison of 8+ aerodynamic designs for faster design decisions.
\resumeItemListEnd

\resumeSubheading
{Handshake AI}{Remote}
{CAD Specialist -- OpenAI-Sponsored AI Evaluation Project}{March 2026 -- Present}
\resumeItemListStart
\item Developed 25+ CAD benchmark tasks and weighted rubrics in SolidWorks, Siemens NX, and AutoCAD to evaluate AI-generated parts, assemblies, and drawings against measurable geometry, dimensioning, and deliverable criteria.
\item Evaluated 40+ developer-created CAD benchmarks and identified 60+ dimensional, manufacturability, feature-quality, and deliverable issues, improving task clarity and scoring consistency before AI model testing.
\resumeItemListEnd

%-----------PROJECTS-----------------
\section{Projects}

\resumeProjectHeading
{Wind Turbine Blade Fluid--Structure Interaction Analysis}
{July 2026}
\resumeProjectItemListStart
\item Integrated CFD and FEA across 6 operating conditions by transferring ANSYS Fluent pressure loads into ANSYS Mechanical, linking aerodynamic loads to blade stress and deformation while validating root radial force within 0.1\% of hand calculations.
\resumeProjectItemListEnd

\resumeProjectHeading
{Unified Aerodynamic Analysis Tool}
{August 2026}
\resumeProjectItemListStart
\item Developed a MATLAB aerodynamic solver combining Prandtl--Glauert-corrected thin-airfoil and shock-expansion theory, then searched 45,000+ configurations to optimize a double-wedge airfoil across subsonic and supersonic regimes.
\resumeProjectItemListEnd

\resumeProjectHeading
{CubeSat Frame \& Attitude Control Design}
{January 2025 -- March 2025}
\resumeProjectItemListStart
\item Reduced CubeSat frame mass by 10\% by iterating structural layout against 5g stiffness, safety-factor, and buckling constraints in ANSYS; developed a MATLAB/Simulink PD controller achieving 3-axis stabilization.
\resumeProjectItemListEnd

%-----------SKILLS-----------------
\section{Skills}

\vspace{+3pt}
\small

\textbf{Simulation:}
ANSYS Fluent (CFD), ANSYS Mechanical (FEA), Simscape Multibody \\

\textbf{3D Modeling:}
SolidWorks (CSWA), CATIA, Inventor, AutoCAD, Fusion 360, Siemens NX \\

\textbf{Systems Engineering:}
MagicDraw/Cameo (SysML), ModelCenter, Simulink \\

\textbf{Programming:}
MATLAB (Certified MATLAB Associate), Python, Git

\end{document}